%!TEX root = ../report.tex
%*******************************************************************************
%                                   Context                                    %
%*******************************************************************************

\chapter{Project context}

%******************************** Section 1.1 *********************************%
\section{Introduction}
  The aim of this chapter is to contextualize the project by giving its background,
  presenting the host company and defining in a later section the main problem as well
  as the way we address it.

%******************************** Section 1.2 *********************************%
\section{Project background}
  This project is about designing and implementing a positive energy worker.
  A Raspberry Pi based system that powers iExec's infrastructure and serves - at the same
  time - as an IoT device to embrace the Edge computing concept.
  It is achieved in the context of the preparations of the end of studies project submitted to
  obtain computer science engineering degree.

%******************************** Section 1.3 *********************************%
\section{Host company}

  \begin{figure}[!h]\centering
    \includegraphics[width=.5\columnwidth]{2-Context/figs/iExec-logo.pdf}
    \caption{iExec logo}
  \end{figure}

  Started in 2016, iExec\cite{iExec} was co-founded by Dr. Gilles Fedak and Pr. Haiwu He.
  
  \subsection{Mission of the company}
    iExec aims at providing decentralized applications running on the blockchain a scalable,
    secure and easy access to the services, data-sets and computing resources they need. This
    technology relies on Ethereum smart contracts and allows the building of a virtual cloud
    infrastructure that provides high-performance computing services on demand.

    iExec leverages a set of research technologies that have been developed at the INRIA and CNRS
    research institutes in the field of Desktop Grid computing. The idea of Desktop Grid
    (aka. Volunteer Computing) is to collect the computer resources that are underutilized on the
    Internet to execute very large parallel applications at the fraction of the cost of a
    traditional supercomputer. iExec relies on XtremWeb-HEP\cite{xtremweb}, a mature, solid, and open-source Desktop
    Grid software which implements all the needed features: fault-tolerance, multi-applications,
    multi-users, hybrid public/private infrastructure, deployment of virtual images, data management,
    security and accountability, and many more.

    iExec is developing a new Proof-of-Contribution\cite{POCO} (PoCo) protocol, that will allow off-chain
    consensus. Thanks to the Proof-of-Contribution, external resource providers will have the usage
    of their resources certified directly in the blockchain.

    iExec aims to deploy a scalable, high-performance, secure and manageable infrastructure sidechain
    that will promote a new form of distributed governance, involving key HPC, big data and cloud
    industry leaders.

  \subsection{Company's crypto-currency}
    iExec is using its ERC20-compliant token to provide standard and secure payments. RLC\cite{RLC} which stands for
    "Run on Lots of Computers", can be securely and easily stored, transferred, traded, divided and used
    to make payments. This widely adopted crypto-currency (87 million RLC are currently in circulation)
    is used to access all iExec's services.

%******************************** Section 1.4 *********************************%
\section{Problematic}
  With the huge increase of human dependency on Information Technology for even daily activities,
  comes the massive consumption of electricity. In 2016, global data centers used roughly 416
  tera-watts (more than 90 billion kilowatt-hours for just U.S. data centers) or about 3\% of the
  total electricity in terms of percentage, which is nearly 40\% more than the entire United Kingdom.
  Predictions have shown that this consumption will double every four years\cite{consumption-prediction}.

  As processing-power demanding technologies like artificial intelligence and blockchain have
  been appearing, the network of data centers that have sprung up in the past decade will spread.
  Moreover, internet-connected devices is changing the entire landscape because IoT is
  projected to exceed 20 billion devices by 2020. Given there are currently 10 billion devices,
  doubling that will require huge increases to our data centers infrastructure, which will
  massively grow our electricity consumption, and that is just adding fuel to the fire.

  There has been some trials to remedy the situation such as using other alternatives to silicon in
  data storage and counting on virtualization to reduce the use of physical machines, but all of that
  still can not keep up against the extreme consumption demand.

%******************************** Section 1.5 *********************************%
\section{Suggested solution}
  Although the existing trials are aiming to make data centers greener, in this project we are taking another
  perspective. The idea is to push computation (or part of it) to the edge of the network instead of sending
  it to the cloud (fog/edge computing principles) and use solar energy to power devices that execute this
  computation. iExec plays a key role in this approach because the execution process will be handled by its
  software, in other words we are creating an iExec's worker that is completely autonomous and energy positive.
  The prototype of this system is a Raspberry Pi based device powered by a solar panel but the concept can be
  applied to a wider range of use cases.

%******************************** Section 1.6 *********************************%
\section{Conclusion}
  After talking about the project's context and background, we presented the host company and discussed the
  problematic we are dealing with as well as the solution to resolve this issue, which is the subject of this
  project.

\clearpage